%%% BEGIN fig:ward-identity
\begin{figure}[htbp]
\centering
\begin{tikzpicture}[fd,
  t2 lab/.style={font=\small},
  t2 eq/.style={font=\large},
  t2 hb/.style={hard blob,ellipse,minimum width=1.3cm,minimum height=0.95cm,inner sep=0pt}]
  % ======================= row (a): one K gluon =======================
  \node[t2 eq] at (-2.35,-0.1) {$\displaystyle\sum_{\text{所有附着}}$};
  \node[t2 hb] (H) at (0,0) {$H$};
  \draw[photon] (-1.25,1.05) -- (H.140);
  \draw[quark] (H.25) -- (1.55,0.85);
  \draw[quark] (-1.5,-1.05) -- (H.214);
  \path[mom={$p$},mom pos=0.4] (-1.5,-1.05) -- (H.214);
  \draw[longit] (-0.55,-1.6) -- (H.264);
  \longitmark{(-0.55,-1.6)}{(H.264)}
  \path[mom'={$k$},mom sep=0.3cm] (-0.55,-1.6) -- (H.264);
  \node[t2 lab,anchor=north] at (-0.55,-1.6) {$a$};
  \node[t2 eq] at (2.35,0) {$=$};
  \begin{scope}[xshift=4.6cm]
    \node[t2 hb] (H0) at (0,0) {$H_0$};
    \draw[photon] (-1.25,1.05) -- (H0.140);
    \draw[quark] (H0.25) -- (1.55,0.85);
    \draw[quark] (-1.5,-1.05) -- (H0.214);
    \path[mom={$p+k$},mom pos=0.4] (-1.5,-1.05) -- (H0.214);
    \draw[wilson] (H0.214) -- ++(-33:2.25) coordinate (W) node[right,t2 lab] {沿 $n_B$};
    \coordinate (Y) at ($(H0.214)!0.52!(W)$);
    \draw[longit] (-0.55,-1.6) -- (Y);
    \longitmark{(-0.55,-1.6)}{(Y)}
    \path[mom'={$k$},mom sep=0.3cm] (-0.55,-1.6) -- (Y);
    \node[t2 lab,anchor=north] at (-0.55,-1.6) {$a$};
  \end{scope}
  \node[font=\small,anchor=west] at (-2.95,-2.4) {(a)\ 阿贝尔情形（任意多个 K 光子），或单个 K 胶子};
  % ======================= row (b): two K gluons, non-abelian =======================
  \begin{scope}[yshift=-4.2cm]
    \node[t2 eq] at (-2.35,-0.1) {$\displaystyle\sum_{\text{所有附着}}$};
    \node[t2 hb] (H) at (0,0) {$H$};
    \draw[photon] (-1.25,1.05) -- (H.140);
    \draw[quark] (H.25) -- (1.55,0.85);
    \draw[quark] (-1.5,-1.05) -- (H.214);
    \draw[longit] (-0.75,-1.62) -- (H.250);  \longitmark{(-0.75,-1.62)}{(H.250)}
    \draw[longit] (0.05,-1.68) -- (H.290);   \longitmark{(0.05,-1.68)}{(H.290)}
    \node[t2 lab,anchor=north] at (-0.75,-1.62) {$k_1$};
    \node[t2 lab,anchor=north] at (0.05,-1.68) {$k_2$};
    \node[t2 eq] at (2.35,0) {$=$};
    \begin{scope}[xshift=4.6cm]
      \node[t2 hb] (H0) at (0,0) {$H_0$};
      \draw[photon] (-1.25,1.05) -- (H0.140);
      \draw[quark] (H0.25) -- (1.55,0.85);
      \draw[quark] (-1.5,-1.05) -- (H0.214);
      \draw[wilson=0.84] (H0.214) -- ++(-33:2.25) coordinate (W) node[right,t2 lab] {沿 $n_B$};
      \draw[longit] (-0.75,-1.62) -- ($(H0.214)!0.36!(W)$);  \longitmark{(-0.75,-1.62)}{($(H0.214)!0.36!(W)$)}
      \draw[longit] (0.05,-1.68) -- ($(H0.214)!0.66!(W)$);   \longitmark{(0.05,-1.68)}{($(H0.214)!0.66!(W)$)}
      \node[t2 lab,anchor=north] at (-0.75,-1.62) {$k_1$};
      \node[t2 lab,anchor=north] at (0.05,-1.68) {$k_2$};
      \node[font=\small,anchor=west] at (1.7,-0.08) {$+\;(k_1\leftrightarrow k_2)$};
    \end{scope}
    \node[t2 eq] at (8.75,0) {$+$};
    \begin{scope}[xshift=10.55cm]
      \node[t2 hb] (Hg) at (0,0) {$H'$};
      \draw[photon] (-1.25,1.05) -- (Hg.140);
      \draw[quark] (Hg.25) -- (1.55,0.85);
      \draw[quark] (-1.5,-1.05) -- (Hg.214);
      \draw[ghost] (-0.85,-1.65) -- (Hg.252);
      \draw[ghost] (Hg.292) -- (0.15,-1.7);
      \node[font=\footnotesize,anchor=north,align=center] at (-0.35,-1.75) {鬼场线\\（鬼场项，示意）};
    \end{scope}
  \end{scope}
  \node[font=\small,anchor=west] at (-2.95,-6.85) {(b)\ 非阿贝尔情形，两个及以上 K 胶子};
\end{tikzpicture}
\caption{图示 Ward 恒等式。蓝色螺旋线是标量（纵向）极化的共线胶子（K 胶子，极化 $\propto k^\mu$，以圆点标记），来自 $A$ 共线子图（未画出）；
$H$ 的其余外线（入射夸克、出射夸克、光子）都取为在壳的物理粒子（共线近似之后）。
(a)~对 K 胶子在 $H$ 内一切可能附着位置（包括非阿贝尔情形中 $H$ 内的胶子线与鬼场线）求和，等于把它附着到一条沿 $n_B$ 方向的 Wilson 线（双线）上、
再乘以不含该胶子的硬部分 $H_0$（其入射动量为 $p+k$）；Wilson 线是未来指向还是过去指向，取决于 K 胶子附着在 $H$ 的末态线还是初态线上。
阿贝尔理论中这一等式对任意多个 K 光子成立；非阿贝尔理论中对单个 K 胶子成立（BRST 对称性）。
(b)~非阿贝尔理论中有两个及以上 K 胶子时，求和给出按路径次序附着在 Wilson 线上的各项（色矩阵 $t^{a_1}$、$t^{a_2}$ 的次序由附着次序决定，两种次序都出现），
一般还伴随含鬼场线（带箭头的点线）的附加项；此处只示意其存在，其精确形式与处理见附录 F。}
\label{fig:ward-identity}
\end{figure}
%%% END fig:ward-identity
